\documentclass[UTF8]{ctexart}
\usepackage{xeCJK}
\usepackage{geometry}
\geometry{left=2cm,right=2cm,top=2.5cm,bottom=2.5cm}
\setlength{\headheight}{12.7pt}

\usepackage{titlesec}
\titleformat{\section}
  {\normalfont\large\bfseries}
  {\thesection}
  {1em}
  {}
\titlespacing*{\section}
  {0pt}
  {3.5ex plus 1ex minus .2ex}
  {2.3ex plus .2ex}

\usepackage{amsmath}
\usepackage{amssymb}
\usepackage{amsthm}
\usepackage{enumitem}
\setlist[enumerate]{itemsep=0pt,topsep=0pt,parsep=0pt,leftmargin=0pt,itemindent=2.5em}
\setlist[itemize]{itemsep=0pt,topsep=0pt,parsep=0pt,leftmargin=0pt,itemindent=2.5em}
\usepackage{fancyhdr}
\pagestyle{fancy}
\fancyhf{}
\usepackage{graphicx}
\usepackage{hyperref}
\usepackage{indentfirst}
\usepackage{lmodern}


\begin{document}
\setlength{\abovedisplayskip}{1pt}
\setlength{\belowdisplayskip}{1pt}

\section{群作用}

% 收到了博士生录取通知书, 大概是此生的最后一张录取通知书.
% 很巧, 正好是结束推研面试的一年之后.
% 2026.07.18

\hypertarget{sec:定义1.1}{}
\noindent\textbf{定义1.1 (群作用)}

给定\href{../01 群、子群与商群/01 半群与群#nameddest=sec:定义2.1}{群}
$G$和集合$X$, 以及映射$\cdot:G\times X\to X$. 如果
\begin{enumerate}
    \item 对任意的$x\in X$, 有$1_G\,x=x$,
    \item 对任意的$g_1,g_2\in G$和$x\in X$, 有$g_1(g_2x)=(g_1g_2)x$,
\end{enumerate}
就称$\cdot$是$G$在$X$上的一个\textbf{群作用}, 记作$\cdot:G\curvearrowright X$.

% 只睡了两个多小时...
% 恭喜西班牙! 早安.
% 2026.07.20

\vspace{10pt}

如果此时有$H<G$, 显然$\cdot\!\restriction_{H\times X}$可以给出$H$在$X$上的群作用.

\vspace{10pt}

\hypertarget{sec:定理1.1}{}
\noindent\textbf{定理1.1}

任给群$G$和集合$X$, 对任意的$\cdot:G\curvearrowright X$, 有群同态
\[
    \overline{\cdot}:G\to S_X,\quad g\mapsto(x\mapsto gx),
\]
并且$\cdot\mapsto\overline{\cdot}$是
$G$在$X$上的群作用与$G$到$S_X$的群同态之间的一一对应.

\vspace{10pt}

\hypertarget{sec:定义1.2}{}
\noindent\textbf{定义1.2}

给定群作用$\cdot:G\curvearrowright X$,
\[
    \mathrm{Ker}\,\overline{\cdot}=\{g\in G\mid\forall x\in X:gx=x\}\\[-3pt]
\]
也称作$\cdot$的核. 如果核是$\{1_G\}$, 就称$\cdot$是\textbf{忠实}的.

\vspace{10pt}

对于一般的群作用$\cdot$\,, $\mathrm{Ker}\,\overline{\cdot}$表示群中
不产生任何作用的元素, 其商群就可以产生忠实的群作用.

\vspace{10pt}

\hypertarget{sec:定理1.2}{}
\noindent\textbf{定理1.2}

对任意的群作用$\cdot:G\curvearrowright X$, 由商群的泛性质, 有群同态
\vspace{3pt}
\[
    \overline{\times}:\,\!^{\displaystyle G}\!
    /_{\displaystyle\mathrm{Ker}\,\overline{\cdot}}\,\to S_X,\\[3pt]
\]
它对应了忠实的群作用$\times:\,\!^{\displaystyle G}\!
/_{\displaystyle\mathrm{Ker}\,\overline{\cdot}}\curvearrowright X$, 并且
对任意的$g\in G$和$x\in X$有$g\mathrm{Ker}\,\overline{\cdot}\times x=gx$.

\vspace{3pt}
\noindent{\kaishu 证明.}

由商群的泛性质, $\mathrm{Ker}\,\overline{\times}=
\,\!^{\displaystyle\mathrm{Ker}\,\overline{\cdot}}\!/_
{\displaystyle\mathrm{Ker}\,\overline{\cdot}}=\{\mathrm{Ker}\,\overline{\cdot}\}$,
从而$\times$忠实, 并且对任意的$g\in G$和$x\in X$有
\vspace{5pt}
\begin{equation*}
    g\mathrm{Ker}\,\overline{\cdot}\times x=
    \overline{\times}(g\mathrm{Ker}\,\overline{\cdot})(x)=
    \overline{\cdot}(g)(x)=gx.
\tag*{\qedsymbol}\end{equation*}

\vspace{15pt}

下面举一些常见的群作用的例子.

\vspace{10pt}

\hypertarget{sec:例1.1}{}
\noindent\textbf{例1.1}

\begin{enumerate}
    \item 任给群$G$, 它在$G$自身上有群作用$(g,g')\mapsto gg'$, 称之为左乘作用.
    \item 任给群$G$和它的正规子群$H$, $G$在$H$
    上有群作用$(g,h)\mapsto ghg^{-1}$, 称之为共轭作用.
    \vspace{3pt}
    \item 任给群$G$和它的子群$H$, $G$在$\,\!^{\displaystyle G}\!/_{\displaystyle H}$
    上有群作用$(g,P)\mapsto gP$, 称之为对陪集的左乘作用.
    \vspace{3pt}
    \item 任给群$G$, $G$在它所有子群组成的集合
    上有群作用$(g,H)\mapsto gHg^{-1}$, 称之为对子群的共轭作用.
\end{enumerate}

\end{document}